\section{AMC 10 10A 2017}
        \begin{enumerate}
        	\item (\textit{AMC 10 10A 2017}) What is the value of $(2(2(2(2(2(2+1)+1)+1)+1)+1)+1)$?


 $\textbf{(A)}\ 70\qquad\textbf{(B)}\ 97\qquad\textbf{(C)}\ 127\qquad\textbf{(D)}\ 159\qquad\textbf{(E)}\ 729$


 
 \textbf{Answer: }C



\textbf{Solution 1}

Notice this is the term $a_6$ in a recursive sequence, defined recursively as $a_1 = 3, a_n = 2a_{n-1} + 1.$ Thus:
 \[\begin{split} a_2 = 3 \cdot 2 + 1 = 7.\\ a_3 = 7 \cdot 2 + 1 = 15.\\ a_4 = 15 \cdot 2 + 1 = 31.\\ a_5 = 31 \cdot 2 + 1 = 63.\\ a_6 = 63 \cdot 2 + 1 = \boxed{\textbf{(C)}\ 127} \end{split}\]

 Minor LaTeX edits by fasterthanlight

 



\textbf{Solution 2}

Starting to compute the inner expressions, we see the results are $1, 3, 7, 15, \ldots$. This is always $1$ less than a power of $2$. The only admissible answer choice by this rule is thus $\boxed{\textbf{(C)}\ 127}$.

 



\textbf{Solution 3}

Working our way from the innermost parenthesis outwards and directly computing, we have $\boxed{\textbf{(C) } 127}$.

 



\textbf{Solution 4}

If you distribute this you get a sum of the powers of $2$.  The largest power of $2$ in the series is $64$, so the sum is $\boxed{\textbf{(C)}\ 127}$.

 


 



\textbf{Solution 5}

$(2(2(2(2(2(2+1)+1)+1)+1)+1)+1)$$=(2(2(2(2(2(3)+1)+1)+1)+1)+1)$$=(2(2(2(2(6+1)+1)+1)+1)+1)$$=(2(2(2(2(7)+1)+1)+1)+1)$$=(2(2(2(14+1)+1)+1)+1)$$=(2(2(2(15)+1)+1)+1)$$=(2(2(30+1)+1)+1)$$=(2(2(31)+1)+1)$$=(2(62+1)+1)$$=(2(63)+1)$$=(126+1)$$=127 \Longrightarrow \boxed{\textbf{(C)}\ 127}$.

 



\textbf{Solution 6 (quickest)}

Notice that $x^6 + x^5 + x^4 + x^3 + x^2 + x + 1 = x (x (x (x (x (x + 1) + 1) + 1) + 1) + 1) + 1$. Substituting $2$ for $x$, we get 
 \[2(2(2(2(2(2+1)+1)+1)+1)+1)+1 = 2^6 + 2^5 + 2^4 + 2^3 + 2^2 + 2 + 1 = 2^7 - 1 \Longrightarrow \boxed{\textbf{(C)}\ 127}\]

 



\textbf{Solution 7 (Fast)}

Notice that there are 5 instances where the sum is multiplied by $2$. That gives us $2^5 = 32$. Separate the addition part into $3$ and the $1$. Multiply the $3$ by $32$, giving $96$. Then notice that the $1$ is multiplied by increasing powers of two; therefore, it is equal to $2^5-1$. Then add these two parts. $96+31=\boxed{\textbf{(C)}\ 127}$

 


	\item (\textit{AMC 10 10A 2017}) Pablo buys popsicles for his friends. The store sells single popsicles for $\$1$ each, $3$-popsicle boxes for $\$2$ each, and $5$-popsicle boxes for $\$3$. What is the greatest number of popsicles that Pablo can buy with $\$8$?


 $\textbf{(A)}\ 8\qquad\textbf{(B)}\ 11\qquad\textbf{(C)}\ 12\qquad\textbf{(D)}\ 13\qquad\textbf{(E)}\ 15$


 
 \textbf{Answer: }D



\textbf{Solution 1}

$\$3$ boxes give us the most popsicles/dollar, so we want to buy as many of those as possible. After buying $2$, we have $\$2$ left. We cannot buy a third $\$3$ box, so we opt for the $\$2$ box instead (since it has a higher popsicles/dollar ratio than the $\$1$ pack). We're now out of money. We bought $5+5+3=13$ popsicles, so the answer is $\boxed{\textbf{(D) }13}$.

 


	\item (\textit{AMC 10 10A 2017}) Tamara has three rows of two $6$-feet by $2$-feet flower beds in her garden. The beds are separated and also surrounded by $1$-foot-wide walkways, as shown on the diagram. What is the total area of the walkways, in square feet?


 
                    \begin{figure}[H]
                        \centering
                        \includegraphics[width=0.5\textwidth]{Problem/AMC_10_10A_2017_3_Problem_Figure_1.png}
                    \end{figure}
                    


 

$\textbf{(A)}\ 72\qquad\textbf{(B)}\ 78\qquad\textbf{(C)}\ 90\qquad\textbf{(D)}\ 120\qquad\textbf{(E)}\ 150$


 
 \textbf{Answer: }B



\textbf{Solution 1}

Finding the area of the shaded walkway can be achieved by computing the total area of Tamara's garden and then subtracting the combined area of her six flower beds.

 Since the width of Tamara's garden contains three margins, the total width is $2\cdot 6+3\cdot 1 = 15$ feet.

 Similarly, the height of Tamara's garden is $3\cdot 2+4\cdot 1 = 10$ feet.

 Therefore, the total area of the garden is $15\cdot 10 =150$ square feet.

 Finally, since the six flower beds each have an area of $2\cdot 6 = 12$ square feet, the area we seek is $150 - 6\cdot 12$, and our answer is $\boxed{\textbf{(B)}\ 78}$

 



\textbf{Solution 2 (probably much less efficient than solution 1)}

The long horizontal walkways have an area of $15 ft^{2}$ each, and there are $4$ of them $\longrightarrow 60 ft^{2}$.

 The short vertical walkways have an area of $2 ft^{2}$ each, and there are $9$ of them $\longrightarrow 18 ft^{2}$.

 Adding the areas together, we have $60 ft^{2} + 18 ft^{2} = 78 ft^{2} \Longrightarrow \boxed{\textbf{(B)}\ 78}$.~JH. L

 


	\item (\textit{AMC 10 10A 2017}) Mia is “helping” her mom pick up $30$ toys that are strewn on the floor. Mia’s mom manages to put $3$ toys into the toy box every $30$ seconds, but each time immediately after those $30$ seconds have elapsed, Mia takes $2$ toys out of the box. How much time, in minutes, will it take Mia and her mom to put all $30$ toys into the box for the first time?


 $\textbf{(A)}\ 13.5\qquad\textbf{(B)}\ 14\qquad\textbf{(C)}\ 14.5\qquad\textbf{(D)}\ 15\qquad\textbf{(E)}\ 15.5$


 
 \textbf{Answer: }B



\textbf{Solution}

Every $30$ seconds, $3$ toys are put in the box and $2$ toys are taken out, so the number of toys in the box increases by $3-2=1$ every $30$ seconds. Then after $27 \times 30 = 810$ seconds (or $13 \frac{1}{2}$ minutes), there are $27$ toys in the box. Mia's mom will then put the remaining $3$ toys into the box after $30$ more seconds, so the total time taken is $27\times 30+30=840$ seconds, or $\boxed{(\textbf{B})\ 14}$ minutes.

 Note: During the last time Mia's mom will complete picking up the $30$ toys (before Mia can take $2$ out) which is the reason that you calculate up to $27$ and then the rest.

 


	\item (\textit{AMC 10 10A 2017}) The sum of two nonzero real numbers is $4$ times their product. What is the sum of the reciprocals of the two numbers?


 $\textbf{(A)}\ 1\qquad\textbf{(B)}\ 2\qquad\textbf{(C)}\ 4\qquad\textbf{(D)}\ 8\qquad\textbf{(E)}\ 12$


 
 \textbf{Answer: }C



\textbf{Solution}

Let the two real numbers be $x,y$. We are given that $x+y=4xy,$ and dividing both sides by $xy$, $\frac{x}{xy}+\frac{y}{xy}=4.$

 
 \[\frac{1}{y}+\frac{1}{x}=\boxed{\textbf{(C) } 4}.\]

 Note: we can easily verify that this is the correct answer; for example, $\left(\frac{1}{2}, \frac{1}{2}\right)$ works, and the sum of their reciprocals is $4$.

 



\textbf{Solution 2}

Instead of using algebra, another approach at this problem would be to notice the fact that one of the nonzero numbers has to be a fraction.See for yourself. And by looking into fractions, we immediately see that $\frac{1}{3}$ and $1$ would fit the rule. $\boxed{\textbf{(C)} 4}.$

 


 



\textbf{Solution 3}

Notice that from the information given above, $x+y=4xy$

 Because the sum of the reciprocals of two numbers is just the sum of the two numbers over the product of the two numbers or $\frac{x+y}{xy}$ 

 We can solve this by substituting $x+y\implies 4xy$.

 Our answer is simply $\frac{4xy}{xy}\implies4$.

 Therefore, the answer is $\boxed{\textbf{(C) } 4}$.

 


 


	\item (\textit{AMC 10 10A 2017}) Ms. Carroll promised that anyone who got all the multiple choice questions right on the upcoming exam would receive an A on the exam. Which of these statements necessarily follows logically?


 $\textbf{(A)}\ \text{If Lewis did not receive an A, then he got all of the multiple choice questions wrong.}\\\textbf{(B)}\ \text{If Lewis did not receive an A, then he got at least one of the multiple choice questions wrong.}\\\textbf{(C)}\ \text{If Lewis got at least one of the multiple choice questions wrong, then he did not receive an A. }\\\textbf{(D)}\ \text{If Lewis received an A, then he got all of the multiple choice questions right.}\\\textbf{(E)}\ \text{If Lewis received an A, then he got at least one of the multiple choice questions right.}$


 
 \textbf{Answer: }B



\textbf{Solution 1 (Using the Contrapositive)}

Rewriting the given statement: "if someone got all the multiple choice questions right on the upcoming exam then he or she would receive an A on the exam."If that someone is Lewis the statement becomes: "if Lewis got all the multiple choice questions right, then he would receive an A on the exam."

 The contrapositive: "If Lewis did \textit{not} receive an A, then he did \textit{not} get all of them right"

 The contrapositive (in other words): "If Lewis did not get an A, then he got at least one of the multiple choice questions wrong".

 B is also equivalent to the contrapositive of the original statement, which implies that it must be true, so the answer is $\boxed{\textbf{(B)}\text{ If Lewis did not receive an A, then he got at least one of the multiple choice questions wrong.}}$.

 
                \begin{itemize}
                    \item Note that answer choice (B) is the contrapositive of the given statement. (That is, it has been negated as well as reversed.) We know that the contrapositive is always true if the given statement is true.
                \end{itemize}


                



\textbf{Solution 2 (Elimination)}

Note: An A is usually 90\%-100\% of the questions correct.

 $\textbf{(A)}$ False. If Lewis did not receive an A, that does not mean he got \textit{all} the multiple-choice questions wrong. For example, he might get 19/20 or 18/20, which still accounts for an A.

 $\textbf{(B)}$ True. If Lewis did not receive an A, then he must have got at least one wrong. Otherwise, Lewis would have gotten an A.

 $\textbf{(C)}$ False. Again, Lewis can get 19/20 or 18/20, which is still an A. 

 $\textbf{(D)}$ False. The above situation can happen.

 $\textbf{(E)}$ False. Lewis can get 17/20 or less but it is not an A.

 Therefore, our answer is $\boxed{\textbf{(B)}}.$

 


 


	\item (\textit{AMC 10 10A 2017}) Jerry and Silvia wanted to go from the southwest corner of a square field to the northeast corner. Jerry walked due east and then due north to reach the goal, but Silvia headed northeast and reached the goal walking in a straight line. Which of the following is closest to how much shorter Silvia's trip was, compared to Jerry's trip?


 $\textbf{(A)}\ 30\%\qquad\textbf{(B)}\ 40\%\qquad\textbf{(C)}\ 50\%\qquad\textbf{(D)}\ 60\%\qquad\textbf{(E)}\ 70\%$


 
 \textbf{Answer: }A



\textbf{Solution}

Let $j$ represent how far Jerry walked, and $s$ represent how far Silvia walked. Since the field is a square, and Jerry walked two sides of it, while Silvia walked the diagonal, we can simply define the side of the square field to be one, and find the distances they walked. Since Jerry walked two sides, $j = 2$Since Silvia walked the diagonal, she walked the hypotenuse of a $45$, $45$, $90$ triangle with leg length $1$. Thus,$s = \sqrt{2} = 1.414...$We can then take $\frac{j-s}{j} \approx \frac{2 - 1.4}{2}=0.3 \implies \boxed{\textbf{(A)}\ 30\%}$

 


 



\textbf{Solution 2}

We can assume that Jerry's path is $100$ units in length. By that logic, Silvia's path is $\frac{100\cdot\sqrt{2}}{2}=50\sqrt{2}$. We can assume that $\sqrt{2}\approx1.4$, so we find that Sylvia's path is $70$ units in length. We can now subtract the two paths, and find out that the answer is $\boxed{\textbf{(A)}\ 30\%}$.~demonhunter

 


	\item (\textit{AMC 10 10A 2017}) At a gathering of $30$ people, there are $20$ people who all know each other and $10$ people who know no one. People who know each other hug, and people who do not know each other shake hands. How many handshakes occur?


 $\textbf{(A)}\ 240\qquad\textbf{(B)}\ 245\qquad\textbf{(C)}\ 290\qquad\textbf{(D)}\ 480\qquad\textbf{(E)}\ 490$


 
 \textbf{Answer: }B



\textbf{Solution 1}

Each one of the ten people has to shake hands with all the $20$ other people they don’t know. So $10\cdot20 = 200$. From there, we calculate how many handshakes occurred between the people who don’t know each other. This is simply counting how many ways to choose two people to shake hands from $10$, or $\binom{10}{2} = 45$. Thus the answer is $200 + 45 = \boxed{\textbf{(B)}\ 245}$.

 



\textbf{Solution 2}

We can also use complementary counting. First of all, $\dbinom{30}{2}=435$ handshakes or hugs occur. Then, if we can find the number of hugs, then we can subtract it from $435$ to find the handshakes. Hugs only happen between the $20$ people who know each other, so there are $\dbinom{20}{2}=190$ hugs. $435-190= \boxed{\textbf{(B)}\ 245}$.

 



\textbf{Solution 3}

We can focus on how many handshakes the $10$ people who don't know anybody get. 

 The first person gets $29$ handshakes with other people not him/herself, the second person gets $28$ handshakes with other people not him/herself and not the first person, ..., and the tenth receives $20$ handshakes with other people not him/herself and not the first, second, ..., ninth person. We can write this as the sum of an arithmetic sequence:

 $\frac{10(20+29)}{2}\implies 5(49)\implies 245.$Therefore, the answer is $\boxed{\textbf{(B)}\ 245}$

 


 



\textbf{Solution 4}

First, we can find out the number of handshakes that the $10$ people who don't know anybody share with the $20$ other people. This is simply $10 \cdot 20 = 200$. Next, we need to find out the number of handshakes that are shared within the $10$ people who don't know anybody. Here, we can use the formula $\frac{n(n-1)}{2}$, where $n$ is the number of people being counted. The reason we divide by $2$ is because $n(n-1)$ counts the case where the $1^{st}$ person shakes hands with the $2^{nd}$ person $and$ the case where the $2^{nd}$ shakes hands with the $1^{st}$ (and these 2 cases are the same). Thus, plugging $n=10$ gives us $\frac{10 \cdot 9}{2} \implies 45$. Adding up the 2 cases gives us $200+45=\boxed{\textbf{(B)}\ 245}$

 


 


	\item (\textit{AMC 10 10A 2017}) Minnie rides on a flat road at $20$ kilometers per hour (kph), downhill at $30$ kph, and uphill at $5$ kph. Penny rides on a flat road at $30$ kph, downhill at $40$ kph, and uphill at $10$ kph. Minnie goes from town $A$ to town $B$, a distance of $10$ km all uphill, then from town $B$ to town $C$, a distance of $15$ km all downhill, and then back to town $A$, a distance of $20$ km on the flat. Penny goes the other way around using the same route. How many more minutes does it take Minnie to complete the $45$-km ride than it takes Penny?


 $\textbf{(A)}\ 45\qquad\textbf{(B)}\ 60\qquad\textbf{(C)}\ 65\qquad\textbf{(D)}\ 90\qquad\textbf{(E)}\ 95$


 
 \textbf{Answer: }C



\textbf{Solution}

The distance from town $A$ to town $B$ is $10$ km uphill, and since Minnie rides uphill at a speed of $5$ kph, it will take her $2$ hours. Next, she will ride from town $B$ to town $C$, a distance of $15$ km all downhill. Since Minnie rides downhill at a speed of $30$ kph, it will take her half an hour. Finally, she rides from town $C$ back to town $A$, a flat distance of $20$ km. Minnie rides on a flat road at $20$ kph, so this will take her $1$ hour. Her entire trip takes her $3.5$ hours.

 Then, Penny will go from town $A$ to town $C$, a flat distance of $20$ km. Since Penny rides on a flat road at $30$ kph, it will take her $\frac{2}{3}$ of an hour. Next Penny will go from town $C$ to town $B$, which is uphill for Penny. Since Penny rides at a speed of $10$ kph uphill, and town $C$ and $B$ are $15$ km apart, it will take her $1.5$ hours. Finally, Penny goes from Town $B$ back to town $A$, a distance of $10$ km downhill. Since Penny rides downhill at $40$ kph, it will only take her $\frac{1}{4}$ of an hour. In total, it takes her $29/12$ hours, which simplifies to $2$ hours and $25$ minutes.Finally, Penny's $2$hr $25$min trip was $\boxed{\textbf{(C)}\ 65}$ minutes less than Minnie's $3$hr $30$min trip.

 


	\item (\textit{AMC 10 10A 2017}) Joy has $30$ thin rods, one each of every integer length from $1$ cm through $30$ cm. She places the rods with lengths $3$ cm, $7$ cm, and $15$ cm on a table. She then wants to choose a fourth rod that she can put with these three to form a quadrilateral with positive area. How many of the remaining rods can she choose as the fourth rod?


 $\textbf{(A)}\ 16\qquad\textbf{(B)}\ 17\qquad\textbf{(C)}\ 18\qquad\textbf{(D)}\ 19\qquad\textbf{(E)}\ 20$


 
 \textbf{Answer: }B



\textbf{Solution 1}

The triangle inequality generalizes to all polygons, so $x < 3+7+15$ and $15<x+3+7$ yields $5<x<25$. Now, we know that there are $19$ numbers between $5$ and $25$ exclusive, but we must subtract $2$ to account for the 2 lengths already used that are between those numbers, which gives $19-2=\boxed{\textbf{(B)}\ 17}$

 



\textbf{Solution 2 (Auxiliary Line)}

If you do not know the fact that the Triangle Inequality applies to all polygons, you can can solve it completely using the standard triangle inequality that you already know by simply splitting the quadrilateral into two triangles.  Let the four sides of our quadrilateral be $3$, $7$, $15$, and $x$. Draw an auxiliary diagonal line $y$ that splits the shape into a small triangle with sides $3$, $7$, and $y$, and a larger triangle with sides $15$, $x$, and $y$.

 Applying the standard triangle inequality to the first triangle with sides $3$, $7$, and $y$ gives us $7 - 3 < y < 7 + 3$, which simplifies to $4 < y < 10$. Note that $y$ is an imaginary line segment slicing through continuous space, meaning it can be any real decimal value in this interval and does not need to be strictly integers.

 Next, we apply the triangle inequality to the second triangle with sides $15$, $x$, and $y$ to determine the boundaries of $x$. For the upper bound, the unknown side must satisfy $x < 15 + y$. As $y$ maximizes and approaches but never reaches $10$, the absolute upper bound approaches $15 + 10 = 25$, meaning $x < 25$. For the lower bound, the side of length $15$ must satisfy $15 < x + y$, which reshapes into $x > 15 - y$. To minimize $x$, we subtract the maximum possible upper limit of $y$ (which approaches $10$), giving $x > 15 - 10 = 5$.

 Combining both boundaries gives us the range $5 < x < 25$. Since $x$ must be an integer, the valid values are $6$ through $24$ inclusive. The total count of integers in this pool is $24 - 6 + 1 = 19$. However, the question states that the rods must be chosen from the remaining pile. The rods of length $7\text{ cm}$ and $15\text{ cm}$ are already placed on the table to build the other sides of the quadrilateral, so they are unavailable to be chosen again. Subtracting these $2$ unavailable choices yields our final total of $19 - 2 = \boxed{\textbf{(B)}\ 17}$.

 


	\item (\textit{AMC 10 10A 2017}) The region consisting of all points in three-dimensional space within $3$ units of line segment $\overline{AB}$ has volume $216\pi$. What is the length $\textit{AB}$?


 $\textbf{(A)}\ 6\qquad\textbf{(B)}\ 12\qquad\textbf{(C)}\ 18\qquad\textbf{(D)}\ 20\qquad\textbf{(E)}\ 24$


 
 \textbf{Answer: }D



\textbf{Solution 1}

In order to solve this problem, we must first visualize what the region looks like. We know that, in a three dimensional space, the region consisting of all points within $3$ units of a point would be a sphere with radius $3$.  However, we need to find the region containing all points within $3$ units of a segment.  It can be seen that our region is a cylinder with two hemispheres/endcaps on either end.  We know the volume of our region, so we set up the following equation (the volume of our cylinder + the volume of our two hemispheres will equal $216 \pi$):

 $\frac{4 \pi }{3} \cdot 3^3+9 \pi x=216 \pi$, where $x$ is equal to the length of our line segment.

 Solving, we find that $x = \boxed{\textbf{(D)}\ 20}$.

 



\textbf{Solution 2}

Because this is just a cylinder and $2$ hemispheres ("half spheres"), and the radius is $3$, the volume of the $2$ hemispheres is $\frac{4(3^3)\pi}{3} = 36 \pi$. Since we also know that the volume of this whole thing is $216 \pi$, we do $216 \pi-36 \pi$ to get $180 \pi$ as the volume of the cylinder. Thus the height is $180 \pi$ divided by the area of the base, or $\frac{180 \pi}{9\pi}=20$, so our answer is $\boxed{\textbf{(D)}\ 20}.$

 \textbf{Diagram for Solution} \href{https://artofproblemsolving.comhttp://i.imgur.com/cwNt293.png}{http://i.imgur.com/cwNt293.png}

 


	\item (\textit{AMC 10 10A 2017}) Let $S$ be a set of points $(x,y)$ in the coordinate plane such that two of the three quantities $3,~x+2,$ and $y-4$ are equal and the third of the three quantities is no greater than this common value. Which of the following is a correct description for $S?$


 $\textbf{(A)}\ \text{a single point} \qquad\textbf{(B)}\ \text{two intersecting lines} \\\qquad\textbf{(C)}\ \text{ three lines whose pairwise intersections are three distinct points} \\\qquad\textbf{(D)}\ \text{a triangle} \qquad\textbf{(E)}\ \text{three rays with a common endpoint}$


 
 \textbf{Answer: }E



\textbf{Solution 1}

If the two equal values are $3$ and $x+2$, then $x=1$. Also, $y-4\le 3$ because $3$ is the common value. Solving for $y$, we get $y \le 7$. Therefore the portion of the line $x=1$ where $y \le 7$ is part of $S$. This is a ray with an endpoint of $(1, 7)$.

 Similar to the process above, we assume that the two equal values are $3$ and $y-4$. Solving the equation $3=y-4$ then $y=7$. Also, $x+2\le 3$ because 3 is the common value. Solving for $x$, we get $x\le1$. Therefore the portion of the line $y=7$ where $x\le 1$ is also part of $S$. This is another ray with the same endpoint as the above ray: $(1, 7)$.

 If $x+2$ and $y-4$ are the two equal values, then $x+2=y-4$. Solving the equation for $y$, we get $y=x+6$. Also $3\le y-4$ because $y-4$ is one way to express the common value. Solving for $y$, we get $y\ge 7$. We also know $3\le x+2$, so $x\ge 1$.Therefore the portion of the line $y=x+6$ where $y\ge 7$ is part of $S$ like the other two rays. The lowest possible value that can be achieved is also $(1, 7)$.

 Since $S$ is made up of three rays with common endpoint $(1, 7)$, the answer is $\boxed{\textbf{(E) }\text{three rays with a common endpoint}}$

 



\textbf{Solution 2 (Graphing \(S\))}

Similar to above, we make three equations.

 \( 3 = x + 2 \)

 \( 3 = y - 4 \)

 \( x + 2 = y - 4 \)

 and solve them in terms of a linear equation

 \( x = 1 \)

 \( y = 7 \)

 \( x + 6 = y \)

 We proceed to graph each of the three equations.

 
                    \begin{figure}[H]
                        \centering
                        \includegraphics[width=0.5\textwidth]{Solution/AMC_10_10A_2017_12_Solution_Figure_1.png}
                    \end{figure}
                     650We now see that no third value is less than the common point. Therefore, we want the overlapping region of \( y \ge 7 \) and \( x \ge 1 \).

 There are three lines in this region that start at a point \((1, 7)\). A line that extends forever starting at any point is called a ray. Because we have three of these lines, particularly \( x = 1 \), \( y = 7 \), and \( x + 6 = y \), we have $\boxed{\textbf{(E) }\text{three rays with a common endpoint}}$.

 


	\item (\textit{AMC 10 10A 2017}) Define a sequence recursively by $F_{0}=0,~F_{1}=1,$ and $F_{n}=$ the remainder when $F_{n-1}+F_{n-2}$ is divided by $3,$ for all $n\geq 2.$ Thus the sequence starts $0,1,1,2,0,2,\ldots$ What is $F_{2017}+F_{2018}+F_{2019}+F_{2020}+F_{2021}+F_{2022}+F_{2023}+F_{2024}?$


 $\textbf{(A)}\ 6\qquad\textbf{(B)}\ 7\qquad\textbf{(C)}\ 8\qquad\textbf{(D)}\ 9\qquad\textbf{(E)}\ 10$


 
 \textbf{Answer: }D



\textbf{Solution}

A pattern starts to emerge as the function is continued. Note that each next term in the sequence is the sum of the previous two terms mod 3. The repeating pattern of 8 terms is $0,1,1,2,0,2,2,1\ldots$ The problem asks for the sum of eight consecutive terms in the sequence, starting from the 2017th. Since the pattern repeats in cycles of 8, we divide 2017 by 8 and we find that the 2017th term is the same as the remainder, the first term. Thus the sum of the sequence from the 2017th to the 2024th term is the same as the sum of the first eight terms. We now just need to find the sum of the first eight numbers, which is just the same as the cycle noted above. Doing this, we get $\boxed{\textbf{(D)}\ 9}$

 


	\item (\textit{AMC 10 10A 2017}) Every week Roger pays for a movie ticket and a soda out of his allowance. Last week, Roger's allowance was $A$ dollars. The cost of his movie ticket was $20\%$ of the difference between $A$ and the cost of his soda, while the cost of his soda was $5\%$ of the difference between $A$ and the cost of his movie ticket. To the nearest whole percent, what fraction of $A$ did Roger pay for his movie ticket and soda?


 $\textbf{(A)}\ 9\%\qquad\textbf{(B)}\ 19\%\qquad\textbf{(C)}\ 22\%\qquad\textbf{(D)}\ 23\%\qquad\textbf{(E)}\ 25\%$


 
 \textbf{Answer: }D



\textbf{Solution}

Let $m$ = cost of movie ticket
Let $s$  = cost of soda

 We can create two equations:

 
 \[m = \frac{1}{5}(A - s)\]
 \[s  = \frac{1}{20}(A - m)\]

 Substituting we get: 


 
 \[m = \frac{1}{5}(A - \frac{1}{20}(A - m))\]which yields:

 \[m = \frac{19}{99}A\]

 Now we can find s and we get:


 
 \[s = \frac{4}{99}A\]

 Since we want to find what fraction of $A$ did Roger pay for his movie ticket and soda, we add $m$ and $s$ to get:

 
 \[\frac{19}{99}A + \frac{4}{99}A \implies \boxed{\textbf{(D)}\ 23\%}\]

 



\textbf{Solution 2}

We have two equations from the problem:$5M=A-S$ and $20S=A-M$If we replace $A$ with $100$ we get a system of equations, and the sum of the values of $M$ and $S$ is the percentage of $A$.Solving, we get $S=\frac{400}{99}$ and $M=\frac{1900}{99}$.Adding, we get $\frac{2300}{99}$, which is closest to $23$ which is $\boxed{\textbf{(D)}\ 23\%}$.

 



\textbf{Solution 3}

Let $m$ be the price of a movie ticket and $s$ be the price of a soda.

 Then, 

 
 \[m=\frac{A-s}{5}\]and
 \[s=\frac{A-m}{20}\]Then, we can turn this into
 \[5m=A-s\]
 \[20s=A-m\]

 Subtracting and getting rid of A, we have $20s-5m=-m+s \rightarrow 19s=4m$. Assuming WLOG (without loss of generality) that $s=4$, $m=19$, thus making a solution for this equation. Substituting this into the 1st equation, we get $A=99$. Hence, $\frac{m+s}{A} = \frac{19+4}{99} \approx \boxed{\textbf{(D)}\ 23\%}$

 


	\item (\textit{AMC 10 10A 2017}) Chloé chooses a real number uniformly at random from the interval $[0, 2017]$. Independently, Laurent chooses a real number uniformly at random from the interval $[0, 4034]$. What is the probability that Laurent's number is greater than Chloé's number? 


 $\textbf{(A)}\ \frac{1}{2}\qquad\textbf{(B)}\ \frac{2}{3}\qquad\textbf{(C)}\ \frac{3}{4}\qquad\textbf{(D)}\ \frac{5}{6}\qquad\textbf{(E)}\ \frac{7}{8}$


 
 \textbf{Answer: }C



\textbf{Solution 1}

We can use geometric probability to solve this.Suppose a point $(x,y)$ lies in the $xy$-plane. Let $x$ be Chloe's number and $y$ be Laurent's number. Then obviously we want $y>x$, which basically gives us a region above a line. We know that Chloe's number is in the interval $[0,2017]$ and Laurent's number is in the interval $[0,4034]$, so we can create a rectangle in the plane, whose length is $2017$ and whose width is $4034$. Drawing it out and dividing into 4 congruent triangles, we see that Laurent's winning area is 3 triangles and Chloe's  is 1 triangle. $\boxed{\textbf{(C)}\ \frac{3}{4}}$.

 



\textbf{Solution 2}

Scale down by $2017$ to get that Chloe picks from $[0,1]$ and Laurent picks from $[0,2]$. There are an infinite number of cases for the number that Chloe picks, but they are all centered around the average of $0.5$. Therefore, Laurent has a winning range of $[X, 2]$, where the average value of $X$ is $0.5$. Thus the average winning length is $2-0.5=1.5$ out of a total length of $2-0=2$. Therefore, the probability is $1.5/2=15/20=\boxed{\frac{3}{4} \space \text{(C)}}.$

 



\textbf{Solution 3 (Fakesolve)}

In total there are $2018\times4034$ ways in which Laurene and Chloe can choose numbers (as same number cannot be chosen by both). If Chloe chooses 2017, then Lauren has 2017 ways to win. If Chloe chooses 2016, Lauren has 2018 ways to win, and so on until Chloe chooses 0, Lauren has 4034 ways to win. Thus the answer is:

 $\frac{\text{number of ways lauren can win}}{\text{number of combinations}}=\frac{2017+2018+2019...+4034}{2018\times4034}$

 Using arithmetic series formula:

 $\frac{\frac{1}{2}(2018)(2017+4034)}{2018\times4034}=\frac{3}{4}$

 $\fbox{C}$ is the correct answer.

 NOTE: Because the problem says real numbers, not integers, this solution may not always work.

 



\textbf{Solution 4 (Intuitive Casework)}

Let Laurene's number \(L \in [0, 4034]\) and Chloé's number \(C \in [0, 2017]\). There is \(\frac{1}{2}\) chance where \(L \in (2017, 4034]\), in which case \(L > C\) regardless of Chloé's number. If \(L \in [0, 2017]\), again by \(\frac{1}{2}\) chance, \(C\) and \(L\) are independently chosen real numbers within the same interval. Then by symmetry, there's equal chance of one being greater than the other; this gives us \(\frac{1}{2} \times \frac{1}{2}\) chance. Here, the chance that they have the same number is very small, so we don't have to account for that situation.

 The answer is \(\frac{1}{2} + \frac{1}{2} \times \frac{1}{2} = \boxed{\textbf{(C)}\ \frac{3}{4}}\)

 



\textbf{Solution 5 (Symmetry / Conditional Probability)}

We can break the problem into two cases based on Laurent's choice:

 
                \begin{itemize}
                    \item **Case 1:** Laurent chooses a number in $(2017, 4034]$. The probability of this happening is $\frac{1}{2}$. Since Chloé's number is at most $2017$, Laurent is guaranteed to win in this case, contributing a probability of $\frac{1}{2} \times 1 = \frac{1}{2}$.
                \end{itemize}


                
                \begin{itemize}
                    \item **Case 2:** Laurent chooses a number in $[0, 2017]$. The probability of this happening is $\frac{1}{2}$. Now, both Laurent and Chloé are choosing independently from the exact same interval $[0, 2017]$. Because they are continuous random variables, the probability of a tie is $0$. By symmetry, Laurent has an equal $\frac{1}{2}$ chance of picking a larger number than Chloé, contributing a probability of $\frac{1}{2} \times \frac{1}{2} = \frac{1}{4}$.
                \end{itemize}


                Summing the probabilities from both cases gives:
 \[\frac{1}{2} + \frac{1}{4} = \boxed{\textbf{(C)}\ \frac{3}{4}}\]

 


	\item (\textit{AMC 10 10A 2017}) There are 10 horses, named Horse 1, Horse 2, $\ldots$, Horse 10. They get their names from how many minutes it takes them to run one lap around a circular race track: Horse $k$ runs one lap in exactly $k$ minutes. At time 0 all the horses are together at the starting point on the track. The horses start running in the same direction, and they keep running around the circular track at their constant speeds. The least time $S>0$, in minutes, at which all 10 horses will again simultaneously be at the starting point is $S=2520$. Let $T>0$ be the least time, in minutes, such that at least 5 of the horses are again at the starting point. What is the sum of the digits of $T$?


 $\textbf{(A)}\ 2\qquad\textbf{(B)}\ 3\qquad\textbf{(C)}\ 4\qquad\textbf{(D)}\ 5\qquad\textbf{(E)}\ 6$


 
 \textbf{Answer: }B



\textbf{Solution 1 (Classical Way)}

If we have horses, $a_1, a_2, \ldots, a_n$, then any number that is a multiple of all those numbers is a time when all horses will meet at the starting point. The least of these numbers is the LCM. To minimize the LCM, we need the smallest primes, and we need to repeat them a lot. By inspection, we find that $\text{LCM}(1,2,3,2\cdot2,2\cdot3) = 12$. Finally, $1+2 = \boxed{\textbf{(B) }3}$.

 



\textbf{Solution 2}

We are trying to find the smallest number that has $5$ one-digit divisors. Therefore we try to find the LCM for smaller digits, such as $1$,$2$, $3$, or $4$. We quickly consider $12$ since it is the smallest number that is the LCM of $1$, $2$, $3$ and $4$. Since $12$ has $5$ single-digit divisors, namely $1$, $2$, $3$, $4$, and $6$, our answer is $1+2 = \boxed{\textbf{(B)}\ 3}$

 



\textbf{Solution 3 (Speedy Guess and Check)}

First, for 5 horses to simultaneously pass the starting line after $T$ seconds, $T$ must be divisible by the amount of seconds it takes each of the 5 horses to pass the starting line, meaning all of the horses must be divisors of $T$, and therefore meaning $T$ must have at least $5$ $1$-digit divisors. Since we want to minimize $T$, we will start by guessing the lowest natural number, $1$. $1$ has only $1$ factor, so it does not work, we now repeat the process for the numbers between $2$ and $12$ (This should not take more than a minute) to get that $12$ is the first number to have $5$ or more single-digit divisors ($1, 2, 3, 4, 6$). The sum of the digits of $12$ is $1+2 = \boxed{\textbf{(B)}\ 3}$, which is our answer.

 Note that this solution is rather fast with this problem because the numbers given in the question are low, this may not always be the case, however, when given higher numbers.

 



\textbf{Solution 4 (Inspection)}

By inspection, $(1, 2, 3, 4, 6)$ yields the lowest answer of $12$ and the sum of the digits is $1+2 \Longrightarrow \boxed{\textbf{(B)}\ 3}$

 


	\item (\textit{AMC 10 10A 2017}) Distinct points $P$, $Q$, $R$, $S$ lie on the circle $x^2+y^2=25$ and have integer coordinates. The distances $PQ$ and $RS$  are irrational numbers. What is the greatest possible value of the ratio $\frac{PQ}{RS}$?


 $\textbf{(A)}\ 3\qquad\textbf{(B)}\ 5\qquad\textbf{(C)}\ 3\sqrt{5}\qquad\textbf{(D)}\ 7\qquad\textbf{(E)}\ 5\sqrt{2}$


 
 \textbf{Answer: }D



\textbf{Solution 1}

Because $P$, $Q$, $R$, and $S$ are lattice points, there are only a few coordinates that actually satisfy the equation. The coordinates are $(\pm 3,\pm 4), (\pm 4, \pm 3), (0,\pm 5),$ and $(\pm 5,0).$ We want to maximize $PQ$ and minimize $RS.$ They also have to be non perfect squares, because they are both irrational. The greatest value of $PQ$ happens when $P$ and $Q$ are almost directly across from each other and are in different quadrants. For example, the endpoints of the segment could be $(-4,3)$ and $(3,-4)$ because the two points are almost across from each other. Another possible pair could be $(-4,3)$ and $(5,0)$. To find out which segment is longer, we have to compare the distances from their endpoints to a diameter (which must be the longest possible segment). The closest diameter would be from $(-4,3)$ to $(4,-3)$. The distance between $(3,-4)$ and $(-4,3)$ is greater than the distance between $(5,0)$ and $(4,-3)$. Therefore, the segment from $(3,-4)$ to $(-4,3)$ is the longest attainable (the other possible coordinates for $P$ and $Q$ are $(4,3)$ and $(-3, -4)$, $(3, 4)$ and $(-4, -3)$, $(-3, 4)$ and $(4, -3)$. The least value of $RS$ is when the two endpoints are in the same quadrant and are very close to each other. This can occur when, for example, $R$ is $(3,4)$ and $S$ is $(4,3).$ They are in the same quadrant and no other point on the circle with integer coordinates is closer to the point $(3,4)$ than $(4,3)$ and vice versa. Using the distance formula, we get that $PQ$ is $\sqrt{98}$ and that $RS$ is $\sqrt{2}.$ $\frac{\sqrt{98}}{\sqrt{2}}=\sqrt{49}=\boxed{\mathrm{\textbf{(D)}}\ 7}$

 



\textbf{Solution 2}

We can look at the option choices. Since we are aiming for the highest possible ratio, let's try using $7$ (though $5 \sqrt{2}$ actually is the highest ratio.) Now, looking at the problem alone, we know that to have the largest ratio possible, we have to let $RS$ be the minimum possible value while at the same time using integer coordinates. Thus, the smallest possible value of $RS$ is $\sqrt{1^{2}+1^{2}} = \sqrt{2}$. Assuming that $\frac{PQ}{RS} = 7$, we plug in $RS = \sqrt{2}$ and solve for $PQ$: $PQ=7\sqrt{2}$. Remember, we don't know if this is possible yet, we are only trying to figure out if it is. But for what values of $x$ and $y$ does $\sqrt{x^{2}+y^{2}}=7\sqrt2$? We see that this can easily be made into a $45-45-90$ triangle. But, instead of substituting $y=x$ into the equation and then using a whole lot of algebra, we can save time and use the little trick, that if in a $45-45-90$ triangle, the two $45$ degree sides have side length $s$, then the hypotenuse is $s\sqrt2$. Using this, we can see that $s=7$, and since our equation does in fact yield a sensible solution, we can be assured that our answer is $\boxed{\mathrm{\textbf{(D)}}\ 7}$.

 Quality Control by fasterthanlight

 (Note by Carrot\_Karen: We tried $7$, but some might be confused why we concluded that it was the answer after verifying without trying the others, like why wasn't option $\textbf{(E)}$ tried? This is because the problem can only have one correct answer, so if we have an option that already works, we can conclude that none of the others work and $\textbf{(D)}$ is the answer.

 comment:(but in other problems, they had different answer choices that worked, so you cant be sure.)

 


	\item (\textit{AMC 10 10A 2017}) Amelia has a coin that lands heads with probability $\frac{1}{3}$, and Blaine has a coin that lands on heads with probability $\frac{2}{5}$. Amelia and Blaine alternately toss their coins until someone gets a head; the first one to get a head wins. All coin tosses are independent. Amelia goes first. The probability that Amelia wins is $\frac{p}{q}$, where $p$ and $q$ are relatively prime positive integers. What is $q-p$?


 $\textbf{(A)}\ 1\qquad\textbf{(B)}\ 2\qquad\textbf{(C)}\ 3\qquad\textbf{(D)}\ 4\qquad\textbf{(E)}\ 5$


 
 \textbf{Answer: }D



\textbf{Solution 1}

Let $P$ be the probability Amelia wins. Note that $P = \text{chance she wins on her first turn} + \text{chance she gets to her turn again}\cdot P$, since if she gets to her turn again, she is back where she started with probability of winning $P$. The chance she wins on her first turn is $\frac{1}{3}$. The chance she makes it to her turn again is a combination of her failing to win the first turn - $\frac{2}{3}$ and Blaine failing to win - $\frac{3}{5}$. Multiplying gives us $\frac{2}{5}$. Thus,
 \[P = \frac{1}{3} + \frac{2}{5}P\]Therefore, $P = \frac{5}{9}$, so the answer is $9-5=\boxed{\textbf{(D)}\ 4}$.

 



\textbf{Solution 2}

Let $P$ be the probability Amelia wins. Note that $P = \text{chance she wins on her first turn} + \text{chance she gets to her second turn}\cdot \frac{1}{3} + \text{chance she gets to her third turn}\cdot \frac{1}{3} \cdots$ This can be represented as an infinite geometric series: 
 \[P=\frac{\frac{1}{3}}{1-\frac{2}{3}\cdot \frac{3}{5}} = \frac{\frac{1}{3}}{1-\frac{2}{5}} = \frac{\frac{1}{3}}{\frac{3}{5}} = \frac{1}{3}\cdot \frac{5}{3} = \frac{5}{9}.\] Therefore, $P = \frac{5}{9}$, so the answer is $9-5 = \boxed{\textbf{(D)}\ 4}.$

 Solution by ktong

 



\textbf{Solution 3}

We can solve this by using 'casework,' the cases being:Case 1: Amelia wins on her first turn.Case 2 Amelia wins on her second turn.and so on.

 The probability of her winning on her first turn is $\dfrac13$. The probability of all the other cases is determined by the probability that Amelia and Blaine all lose until Amelia's turn on which she is supposed to win. So, the total probability of Amelia winning is:
 \[\dfrac{1}{3}+\left(\dfrac{2}{3}\cdot\dfrac{3}{5}\right)\cdot\dfrac{1}{3}+\left(\dfrac{2}{3}\cdot\dfrac{3}{5}\right)^2\cdot\dfrac{1}{3}+\cdots.\]Factoring out $\dfrac13$ we get a geometric series:
 \[\dfrac{1}{3}\left(1+\dfrac{2}{5}+\left(\dfrac{2}{5}\right)^2+\cdots\right) = \dfrac{1}{3}\cdot\dfrac{1}{3/5} = \boxed{\dfrac59}.\]

 Extracting the desired result, we get $9-5 = \boxed{\textbf{(D)} \ 4}$.

 -ConfidentKoala4

 


	\item (\textit{AMC 10 10A 2017}) Alice refuses to sit next to either Bob or Carla. Derek refuses to sit next to Eric. How many ways are there for the five of them to sit in a row of 5 chairs under these conditions?


 $\textbf{(A)}\ 12\qquad\textbf{(B)}\ 16\qquad\textbf{(C)}\ 28\qquad\textbf{(D)}\ 32\qquad\textbf{(E)}\ 40$


 
 \textbf{Answer: }C



\textbf{Solution 1: Casework}

Let Alice be A, Bob be B, Carla be C, Derek be D, and Eric be E.We can split this problem up into two cases:

 $\textbf{Case 1: }$ A sits on an edge seat.

 Since B and C can't sit next to A, that must mean either D or E sits next to A. After we pick either D or E, then either B or C must sit next to D/E. Then, we can arrange the two remaining people in two ways. Since there are two different edge seats that A can sit in, there are a total of $2 \cdot 2 \cdot 2 \cdot 2 = 16$.

 $\textbf{Case 2: }$ A does not sit in an edge seat.

 Still, the only two people that can sit next to A are either D or E, and there are two ways to permute them, and this also handles the restriction that D can't sit next to E. Then, there are two ways to arrange B and C, the remaining people. However, there are three initial seats that A can sit in, so there are $3 \cdot 2 \cdot 2 = 12$ seatings in this case.

 Adding up all the seatings, we have $16+12 = \boxed{\textbf{(C) } 28}$.

 Minor clarity edits by SwordAxe

 



\textbf{Solution 2}

Label the seats (from left to right) $1$ through $5$. The number of ways to seat Derek and Eric in the five seats with no restrictions is $5 \cdot 4=20$. The number of ways to seat Derek and Eric such that they sit next to each other is $8$ (we can treat Derek and Eric as a "block". There are four ways to seat this "block", and two ways to permute Derek and Eric, for a total of $4\cdot 2=8$), so the

 


	\item (\textit{AMC 10 10A 2017}) Let $S(n)$ equal the sum of the digits of positive integer $n$. For example, $S(1507) = 13$. For a particular positive integer $n$, $S(n) = 1274$. Which of the following could be the value of $S(n+1)$?


 $\textbf{(A)}\ 1 \qquad\textbf{(B)}\ 3\qquad\textbf{(C)}\ 12\qquad\textbf{(D)}\ 1239\qquad\textbf{(E)}\ 1265$


 
 \textbf{Answer: }D



\textbf{Solution 1}

Note that $n \equiv S(n) \pmod{9}$. This can be seen from the fact that $\sum_{k=0}^{n}10^{k}a_k \equiv \sum_{k=0}^{n}a_k \pmod{9}$. Thus, if $S(n) = 1274$, then $n \equiv 5 \pmod{9}$, and thus $n+1 \equiv S(n+1) \equiv 6 \pmod{9}$. The only answer choice that satisfies $n+1 \equiv 6 \pmod{9}$ is $\boxed{\textbf{(D)} 1239}$.

 



\textbf{Solution 2 (Sol. 1 but more detail)}

We disregard the case that the sum of $S(N)$ is $0$ because that would indicate that the number must be $0$ itself, which is not allowed. Therefore, every positive integer $N \ge 1$ can be denoted in the modular 9 format with $N \bmod 9$. Thus we have $S(N) \bmod 9 \). In this case, we have \( S(n) = 1274$, and therefore we have $S(n) \bmod 9 \equiv 1274 \bmod 9 \equiv 5 \bmod 9$. Thus, \( S(n+1) \) must be congruent to $6 \bmod 9$, and the only value that satisfies this is $\boxed{\textbf{(D)}\ 1239}$.

 



\textbf{Solution 3}

One divisibility rule that we can use for this problem is that a multiple of $9$ will always have its digits sum to a multiple of $9$. We can find out that the least number of digits the number $n$ has is $142$, with $141$ $9$'s and $1$ $5$, assuming the rule above. No matter what arrangement or different digits we use, the divisibility rule stays the same. To make the problem simpler, we can just use the $141$ $9$'s and $1$ $5$. By randomly mixing the digits up, we are likely to get: $9999$...$9995999$...$9999$. By adding $1$ to this number, we get: $9999$...$9996000$...$0000$.Knowing that $n+1$ is divisible by $9$ when $6$, we can subtract $6$ from every available choice, and see if the number is divisible by $9$ afterwards.After subtracting $6$ from every number, we can conclude that $1233$ (originally $1239$) is the only number divisible by $9$. So our answer is $\boxed{\textbf{(D)}\ 1239}$.

 



\textbf{Solution 4}

The number $n$ can be viewed as having some unique digits in the front, following by a certain number of nines. We can then evaluate each potential answer choice.

 If $A$ is correct, then $n$ must be some number $99999999...9$, because when we add one to $99999999...9$ we get $10000000...00$. Thus, if $1$ is the correct answer, then the equation $9x=1274$ must have an integer solution (i.e. $1274$ must be divisible by $9$). But since it does not, $1$ is not the correct answer.

 If $B$ is correct, then $n$ must be some number $29999999...9$, because when we add one to $29999999...9$, we get $30000000...00$. Thus, if $3$ is the correct answer, then the equation $2+9x=1274$ must have an integer solution. But since it does not, $3$ is not the correct answer.

 Based on what we have done for evaluating the previous two answer choices, we can create an equation we can use to evaluate the final three possibilities. Notice that if $S(n+1)=N$, then $n$ must be a number whose initial digits sum to $N-1$, and whose other, terminating digits, are all $9$. Thus, we can evaluate the three final possibilities by seeing if the equation $(N-1)+9x=1274$ has an integer solution.

 The equation does not have an integer solution for $N=12$, so $C$ is not correct. However, the equation does have an integer solution for $N=1239$ ($x=4$), so $\boxed{\textbf{(D)} ~1239}$ is the answer.

 



\textbf{Solution 5 (Intuition)}

If adding $1$ to $n$ does not carry any of its digits, then $S(n+1)=S(n)+1$ (ex: $25+1=26$. Sum of digits $7 \rightarrow 8$). But since no answer choice is $1275$, that means $n$ has some amount of $9$'s from right to left. 

 When $n+1$, some $9$'s will bump to 0, not affected its$\pmod 9$. But the first non-9 digit (from right to left) will be bumped up by 1. So $S(n) + 1 \pmod {9} \equiv S(n+1) \pmod{9}$. For example, $34999+1=35000$, and the sum of digits $7+27 \rightarrow 8+0$.

 Since $S(n) \equiv 5 \pmod{9}$, that means $S(n+1) \equiv 6 \pmod{9}$. The only answer choice that meets this requirement is $\boxed{\textbf{(D) } 1239}.$

 



\textbf{Solution 6 (Answer choices and luck)}

We note quickly from the given value of $S(n)$ that the form of n is likely going to be of the form $999999999...X9...9999$  where X is a natural number between 1 and 9 inclusive

 This allows us to create the following equation

 
$S(n+1)=1275 - 9y= A$

 where $A$ is an answer choice and $y$ is any natural number

 Plugging in answer choices for $A$ we quickly see that $\boxed{\textbf{(D) } 1239}.$ is our answer

 $\textbf{Note}$: we do have a second possible form for $n$ which is 

 $99999...993$ 

 and thus $n+1 = 99999...994$ and $S(n+1)=1273$

 but since $1273$ is not an answer choice we can disregard this case

 



\textbf{Solution 7 (Answer choices but what's luck?)}

This is a very complex solution so I recommend you read the previous ones before this.

 Looking at the answer choices we first take A.

 We see that A cannot work as if it was 1, that means that the number must be 1 followed by \( n \) amounts of zero's, and this cannot work as \( S(n) \) itself is not divisible by 9.

 We look at 3 and notice that we have two cases. Either it is \( 1 + 2 + n \) zero's or it is 3 followed by \( n \) zero's. Case 2 is obviously false as subtracting 2 still doesn't make \( S(n) \) divisible by 9, and \( 1 + 2 + n \) zero's also doesn't work as even when we subtract 3 from the total, \( S(n) \) proceeds to still not be divisible by 9.

 We look at 12 an immediately notice this can't be true as it needs to be divisible by 3 and 4 and the number itself cannot be both divisible by 3 and by 2 as when adding 1 to \( S(n) \) gives 1275, which is odd.

 We look at 1239 and see that the sum of the digits is exactly 1 greater than \( S(n) \), making it a possibility.

 We see that 1265 has a sum that is exactly the same as \( S(n) \). This means that the subaverage of the number lies between an interval of 9. But the sum \( S(n) \) itself is not divisible by 9, proving that the subaverage is false.

 Our only answer is $\boxed{\textbf{(D) } 1239}.$

 



\textbf{Solution 8}

If $n$ ends in a number that is not $9$, then $S(n+1)=1275$. As this is not one of the answer choices, we can conclude that $n$ must end with $9$. 

 Since $1274$ is not a multiple of $9$, we know that $n$ cannot be entirely made of $9$'s. Therefore, we must have, somewhere along $n$, a digit(we can call it $a$) that isn't $9$, followed by at least one $9$ (since $n$ ends in a $9$). Let the number of $9$'s that occur after $a$ be $x$. 

 When we add $1$ to $n$ to make $n+1$, all of the $9$'s that come after $a$ become $0$ and $a$ becomes the digit $a+1$. Therefore, since we lose the value that the $9$ x's contributed and gain $1$,  $S(n+1)=S(n)-9x+1$. We know that $x$ is an integer, and $S(n)-S(n+1)=9x-1$. Furthermore, we know that $S(n)=1274$, so $1274-S(n+1)$ is one less than a multiple of $9$. 

 The only answer choice that satisfies this is  $\boxed{\textbf{(D) } 1239}.$

 


	\item (\textit{AMC 10 10A 2017}) A square with side length $x$ is inscribed in a right triangle with sides of length $3$, $4$, and $5$ so that one vertex of the square coincides with the right-angle vertex of the triangle. A square with side length $y$ is inscribed in another right triangle with sides of length $3$, $4$, and $5$ so that one side of the square lies on the hypotenuse of the triangle. What is $\frac{x}{y}$?


 $\textbf{(A) } \dfrac{12}{13} \qquad \textbf{(B) } \dfrac{35}{37} \qquad \textbf{(C) } 1 \qquad \textbf{(D) } \dfrac{37}{35} \qquad \textbf{(E) } \dfrac{13}{12}$


 
 \textbf{Answer: }D



\textbf{Solution 1}

Analyze the first right triangle.

 
                    \begin{figure}[H]
                        \centering
                        \includegraphics[width=0.5\textwidth]{Solution/AMC_10_10A_2017_21_Solution_Figure_1.png}
                    \end{figure}
                    

 Note that $\triangle ABC$ and $\triangle FBE$ are similar, so $\frac{BF}{FE} = \frac{AB}{AC}$. This can be written as $\frac{4-x}{x}=\frac{4}{3}$. Solving, $x = \frac{12}{7}$.

 Now we analyze the second triangle.

 

                    \begin{figure}[H]
                        \centering
                        \includegraphics[width=0.5\textwidth]{Solution/AMC_10_10A_2017_21_Solution_Figure_2.png}
                    \end{figure}
                    

 Similarly, $\triangle A'B'C'$ and $\triangle RB'Q$ are similar, so $RB' = \frac{4}{3}y$, and $C'S = \frac{3}{4}y$. Thus, $C'B' = C'S + SR + RB' = \frac{4}{3}y + y + \frac{3}{4}y = 5$. Solving for $y$, we get $y = \frac{60}{37}$. Thus, $\frac{x}{y} = \boxed{\textbf{(D)}\:\frac{37}{35}}$.

 


 



\textbf{Solution 2 (Full Coordinate Bash for x and y)}

Set the right-angle vertex of the triangle as $(0,0)$. Notice that the hypotenuse of the triangle can be described by the line equation $y = 3 - \frac{3}{4}x$, which simplifies to $3x + 4y - 12 = 0$. 

 To find $x$, the top-right vertex of the first square lies on this hypotenuse and can be described by the line $y=x$. Hence, we may solve for $x$ by solving $3x + 4x = 12$, which yields $x = \frac{12}{7}$.

 To find $y$, let the two corners of the second square that touch the sides of the triangle be $Q(x_Q, 0)$ and $T(0, y_T)$. Pushing these points straight out by a distance of $y$ gives the other two corners on the hypotenuse:
 \[R = \left(x_Q + \frac{3}{5}y, \frac{4}{5}y\right)\]
 \[S = \left(\frac{3}{5}y, y_T + \frac{4}{5}y\right)\]

 Since $R$ and $S$ lie on the hypotenuse line $3x + 4y = 12$, we can plug them in to find $x_Q$ and $y_T$:
 \[3\left(x_Q + \frac{3}{5}y\right) + 4\left(\frac{4}{5}y\right) = 12 \implies x_Q = \frac{12-5y}{3}\]
 \[3\left(\frac{3}{5}y\right) + 4\left(y_T + \frac{4}{5}y\right) = 12 \implies y_T = \frac{12-5y}{4}\]

 The distance between the two points $Q$ and $T$ is just the side of the square $y$. Using the distance formula:
 \[\left(\frac{12-5y}{3}\right)^2 + \left(\frac{12-5y}{4}\right)^2 = y^2\]
 \[(12-5y)^2 \left(\frac{1}{9} + \frac{1}{16}\right) = y^2\]
 \[(12-5y)^2 \left(\frac{25}{144}\right) = y^2\]

 Taking the square root of both sides gives:
 \[\frac{5}{12}(12-5y) = y \implies 60 - 25y = 12y \implies 37y = 60 \implies y = \frac{60}{37}\]

 Thus, $\frac{x}{y} = \frac{12/7}{60/37} = \boxed{\textbf{(D)}\:\frac{37}{35}}$.

 



\textbf{Solution 3 (Alternate solution in finding x)}

Set the right-angle vertex of the triangle as $(0,0)$. Notice that the hypotenuse of the triangle, as depicted in solution one, can be described by $y = 3 - \frac{3}{4}x$, while $AE$ can be describe by $y=x$. Hence, we may solve for $x$ by solving $3- \frac{3}{4}x = x$, which yields $\frac{12}{7}$.

 Proceed by finding the value of y via the method described in solution 1, and we will get $y = \frac{60}{37}$. Thus, $\frac{x}{y} = \boxed{\textbf{(D)}\:\frac{37}{35}}$.

 \textbf{Note} In general, if the legs were $a$ and $b$, we have that 
 \[x= \frac{ab}{a+b}, y=\frac{ab\sqrt{a^2+b^2}}{a^2+ab+b^2}.\] This can be verified by plugging in $a=3$ and $b=4$.~anduran

 


	\item (\textit{AMC 10 10A 2017}) Sides $\overline{AB}$ and $\overline{AC}$ of equilateral triangle $ABC$ are tangent to a circle at points $B$ and $C$ respectively.  What fraction of the area of $\triangle ABC$ lies outside the circle?


 $\textbf{(A)}\ \frac{4\sqrt{3}\pi}{27}-\frac{1}{3}\qquad\textbf{(B)}\ \frac{\sqrt{3}}{2}-\frac{\pi}{8}\qquad\textbf{(C)}\ \frac{1}{2}\qquad\textbf{(D)}\ \sqrt{3}-\frac{2\sqrt{3}\pi}{9}\qquad\textbf{(E)}\ \frac{4}{3}-\frac{4\sqrt{3}\pi}{27}$


 
 \textbf{Answer: }E



\textbf{Solution 1}


                    \begin{figure}[H]
                        \centering
                        \includegraphics[width=0.5\textwidth]{Solution/AMC_10_10A_2017_22_Solution_Figure_1.png}
                    \end{figure}
                    Let the radius of the circle be $r$, and let its center be $O$. Since $\overline{AB}$ and $\overline{AC}$ are tangent to circle $O$, then $\angle OBA = \angle OCA = 90^{\circ}$, so $\angle BOC = 120^{\circ}$. Therefore, since $\overline{OB}$ and $\overline{OC}$ are equal to $r$, then (pick your favorite method) $\overline{BC} = r\sqrt{3}$. The area of the equilateral triangle is $\frac{(r\sqrt{3})^2 \sqrt{3}}4 = \frac{3r^2 \sqrt{3}}4$, and the area of the sector we are subtracting from it is $\frac 13 \pi r^2 - \frac 12 \cdot r\sqrt{3} \cdot \frac{r}2 = \frac{\pi r^2}3 -\frac{r^2 \sqrt{3}}4$. The area outside of the circle is $\frac{3r^2 \sqrt{3}}4-\left(\frac{\pi r^2}3 -\frac{r^2 \sqrt{3}}4\right) = r^2 \sqrt{3} - \frac{\pi r^2}3$. Therefore, the \href{https://artofproblemsolving.com/wiki/index.php?title=Answer}{answer} is 
 \[\frac{r^2 \sqrt{3} - \frac{\pi r^2}3}{\frac{3r^2 \sqrt{3}}4} = \boxed{\textbf{(E) } \frac 43 - \frac{4\sqrt 3 \pi}{27}}\]

 \textbf{Note} The sector angle is $120$ because $\angle B$ and $\angle C$ are both 90 degrees meaning $\angle B + \angle C = 180^\circ$, so $ABCO$ is cyclic. Thus, \href{https://artofproblemsolving.com/wiki/index.php?title=The&action=edit&redlink=1}{the} angle is $180-60=120^\circ$

 Alternately, $\angle ABC$ is $60^\circ$ and $\angle ABO$ is $90^\circ$, making $\angle CBO$ $30^\circ$. Symmetry allows us to use the same argument to get $\angle BCO = 30^\circ$. Since the interior angles of $\triangle BCO$ must sum to $180^\circ$, that leaves $120^\circ$ for \href{https://artofproblemsolving.com/wiki/index.php?title=Central\_angle}{central angle} $\angle BOC$.

 —\href{https://artofproblemsolving.com/wiki/index.php?title=User:Wescarroll&action=edit&redlink=1}{wescarroll}

 


 \begin{verbatim}
The 90 degrees comes from the tangency.
\end{verbatim}

\textbf{Multiple Choice Shortcut} Assuming WLoG that the equilateral triangle's side length $s$ and therefore area $A$ are algebraic ("$\pi$-free"):

 The "crust" is a circle sector minus a triangle, so its area is $a \pi - b$, where $a$ and $b$ are algebraic. Thus the answer is $(A - (a \pi -b))/A = (A+b)/A - a\pi/A$.

 Once you see that $s$ is $\sqrt{3} \times$ the circle's radius, and that the circle's 30°-30°-120° triangle is two halves of an equilateral triangle, infer that the smaller circle-sector triangle's area is $A/(\sqrt 3)^2 = A/3$, and so the algebraic part of the answer $(A+b)/A = (A - (-A/3)) /A = 4/3$.

 The transcendental ("$\pi$") part of the answer is $-a \pi /A$, and since $a$ and $A$ are algebraic,  $\textbf{(E)}$ is the only compatible answer choice.

 



\textbf{Solution 2}


                    \begin{figure}[H]
                        \centering
                        \includegraphics[width=0.5\textwidth]{Solution/AMC_10_10A_2017_22_Solution_Figure_2.png}
                    \end{figure}
                    (same diagram as Solution 1)

 Without the Loss of Generality, let the side length of the triangle be $1$. 

 Then, the area of the triangle is $\frac{\sqrt{3}}{4}$. We are looking for the area of the portion inside the triangle but outside the circle divided by the area of the triangle. Since $m\angle ABO = 90^{\circ}=m\angle ACO = 90^{\circ}$, and $m\angle ABC = m\angle ACB = 60^{\circ}$, we know $m\angle OBC=m\angle OCB=30^{\circ}$, and $m\angle BOC = 120^{\circ}$. Drop an angle bisector of $O$ onto $BC$, call the point of intersection $D$. By SAS congruence, $\triangle BDO \cong \triangle CDO$, by CPCTC (Congruent Parts of Congruent Triangles are Congruent) $BD \cong DC$ and they both measure $\frac{1}{2}$. By 30-60-90 triangle, $OC = BO = \frac{\sqrt{3}}{3}$. The area of the sector bounded by arc BC is one-third the area of circle O, whose area is $\left(\frac{\sqrt{3}}{3}\right)^{2}\pi=\frac{1}{3}\pi$. Therefore, the area of the sector bounded by arc BC is $\frac{1}{3}\cdot\frac{1}{3}\pi=\frac{\pi}{9}$. 

 We are nearly there. By 30-60-90 triangle, we know $DO = \frac{\sqrt{3}}{6}$, so the area of $\triangle BOC$ is $\frac{1\cdot\frac{\sqrt{3}}{6}}{2}=\frac{\sqrt{3}}{12}$. The area of the region inside both the triangle and circle is the area of the sector bounded by arc BC minus the area of $\triangle BOC$: $\frac{\pi}{9}-\frac{\sqrt{3}}{12}$. The area of the region outside of the circle but inside the triangle is $\frac{\sqrt{3}}{4}-\left(\frac{\pi}{9}-\frac{\sqrt{3}}{12}\right)=\frac{\sqrt{3}}{4}-\frac{\pi}{9}+\frac{\sqrt{3}}{12}=\frac{\sqrt{3}}{3}-\frac{\pi}{9}$ and the ratio is $\frac{\frac{\sqrt{3}}{3}-\frac{\pi}{9}}{\frac{\sqrt{3}}{4}}=\left(\frac{\sqrt{3}}{3}-\frac{\pi}{9}\right)\cdot\frac{4}{\sqrt{3}}=\frac{4\sqrt{3}}{3\sqrt{3}}-\frac{4\pi}{9\sqrt{3}}=\frac{4}{3}-\frac{4\sqrt{3}\pi}{27} \Longrightarrow \boxed{\textbf{(E) } \frac 43 - \frac{4\sqrt 3 \pi}{27}}$.

 


	\item (\textit{AMC 10 10A 2017}) How many triangles with positive area have all their vertices at points $(i,j)$ in the coordinate plane, where $i$ and $j$ are integers between $1$ and $5$, inclusive?


 $\textbf{(A)}\ 2128 \qquad\textbf{(B)}\ 2148 \qquad\textbf{(C)}\ 2160 \qquad\textbf{(D)}\ 2200 \qquad\textbf{(E)}\ 2300$


 
 \textbf{Answer: }B



\textbf{Solution 1}

We can solve this by finding all the combinations, then subtracting the ones that are on the same line. There are $25$ points in all, from $(1,1)$ to $(5,5)$, so $\dbinom{25}3$ is $\frac{25\cdot 24\cdot 23}{3\cdot 2 \cdot 1}$, which simplifies to $2300$. Now we count the ones that are on the same line. We see that any three points chosen from $(1,1)$ and $(1,5)$ would be on the same line, so $\dbinom53$ is $10$, and there are $5$ rows, $5$ columns, and $2$ long diagonals, so that results in $120$. We can also count the ones with $4$ on a diagonal. That is $\dbinom43$, which is 4, and there are $4$ of those diagonals, so that results in $16$. We can count the ones with only $3$ on a diagonal, and there are $4$ diagonals like that, so that results in $4$. We can also count the ones with a slope of $\frac12$, $2$, $-\frac12$, or $-2$, with $3$ points in each. Note that there are $3$ such lines, for each slope, present in the grid. In total, this results in $12$. Finally, we subtract all the ones in a line from $2300$, so we have $2300-120-16-4-12=\boxed{(\textbf{B})\ 2148}$

 



\textbf{Solution 2}

There are $5 \times 5 = 25$ total points in all. So, there are $\dbinom{25}3 = 2300$ ways to choose the three vertices for the triangle. However, there are some cases where they 3 points chosen results in a straight line.

 There are $10 \times 10 = 100$ cases where the 3 points chosen make up a vertical or horizontal line.

 There are $2\left(1+\dbinom{4}3+\dbinom{5}3+\dbinom{4}3+1\right)=40$ cases where the 3 points all land on the diagonals of the square.

 There are $3 \times 4=12$ ways where the 3 points make the a slope of $\frac{1}{2}$, $-\frac{1}{2}$, $2$, and $-2$.

 Hence, there are $100+40+12=152$ cases where the chosen 3 points make a line. The answer would be $2300-152=\boxed{(\textbf{B})\ 2148}$

 



\textbf{Solution 3 (Counting Triplets by Line Length)}

We can solve this by grouping the straight lines based on how many points are actually on them. If a line has $k$ points on it, the number of bad triplets we can pull from it is just $\dbinom{k}{3}$. 

 First, the total number of ways to pick any 3 points out of the 25 coordinates is $\dbinom{25}{3} = \frac{25 \cdot 24 \cdot 23}{3 \cdot 2 \cdot 1} = 2300$. 

 Now, we count the lines based on their point counts to find all the bad triplets:

 
                \begin{itemize}
                    \item \textbf{Lines with exactly 5 points:} These are the full-length lines across the grid. We have $5$ rows, $5$ columns, and $2$ long main diagonals. This gives $5 + 5 + 2 = 12$ lines total.
                \end{itemize}


                Triplets from this group: $12 \times \dbinom{5}{3} = 12 \times 10 = 120$.

 
                \begin{itemize}
                    \item \textbf{Lines with exactly 4 points:} These are the diagonals directly next to the main ones. There are 2 sloping up and 2 sloping down, which gives $4$ lines total.
                \end{itemize}


                Triplets from this group: $4 \times \dbinom{4}{3} = 4 \times 4 = 16$.

 
                \begin{itemize}
                    \item \textbf{Lines with exactly 3 points:} This group captures the remaining shorter lines and the hidden slopes:
                \end{itemize}


                
            \begin{enumerate}
                \item The tiny corner diagonals ($2$ sloping up, $2$ sloping down) $\rightarrow$ 4 lines.\item The "knight-move" lines with slopes of $\pm 2$ and $\pm \frac{1}{2}$. For each of these 4 unique slopes, there are exactly 3 parallel lines that hit exactly 3 points within the grid boundaries $\rightarrow$ $4 \times 3 = 12$ lines.
            \end{enumerate}
            This gives a total of $4 + 12 = 16$ lines with exactly 3 points. Triplets from this group: $16 \times \dbinom{3}{3} = 16 \times 1 = 16$.

 Adding up all the bad triplets gives: $120 + 16 + 16 = 152$.

 Subtracting these from the total pool gives the number of triangles with positive area: $2300 - 152 = \boxed{(\textbf{B})\ 2148}$.

 


	\item (\textit{AMC 10 10A 2017}) For certain real numbers $a$, $b$, and $c$, the polynomial 
 \[g(x) = x^3 + ax^2 + x + 10\]has three distinct roots, and each root of $g(x)$ is also a root of the polynomial 
 \[f(x) = x^4 + x^3 + bx^2 + 100x + c.\]What is $f(1)$?


 $\textbf{(A)}\ -9009\qquad\textbf{(B)}\ -8008\qquad\textbf{(C)}\ -7007\qquad\textbf{(D)}\ -6006\qquad\textbf{(E)}\ -5005$


 
 \textbf{Answer: }C



\textbf{Solution 1}

$f(x)$ must have four roots, three of which are roots of $g(x)$. Using the fact that every polynomial has a unique factorization into its roots, and since the leading coefficient of $f(x)$ and $g(x)$ are the same, we know that

 
 \[f(x)=g(x)(x+p)\]

 where $-p\in\mathbb{R}$ is the fourth root of $f(x)$. (Using $(x+p) = (x-r))$ instead of $(x-r)$ makes the following computations less messy.)Substituting $g(x)$ and expanding, we find that

 
 \begin{align*}f(x)&=(x^3+ax^2+x+10)(x+p)\\ &=x^4+(a+p)x^3+(1+ap)x^2+(10+p)x+10p.\end{align*}


 Comparing coefficients with $f(x)$, we see that

 
 \begin{align*} a+p&=1\\ 1+ap=b\\ 10+p&=100\\ 10p&=c.\\ \end{align*}


 
Let's solve for $a,b,c,$ and $p$. Since $10+p=100$, $p=90$.

 Since $a+p=1$, $a=-89$. 

 (Solution 1.1 branches from here and takes a shortcut.)

 $c=(10)(90)=900$.

 Then, since $b=1+ap$, $b=-8009$. Thus,

 
 \[f(x)=x^4+x^3-8009x^2+100x+900.\]

 (Solution 1.2 branches from here and takes another shortcut)

 Taking $f(1)$, we find that

 
 \begin{align*} f(1)&=1^4+1^3-8009(1)^2+100(1)+900\\ &=1+1-8009+100+900\\ &=\boxed{\bold{(C)}\, -7007}.\\ \end{align*}


 



\textbf{Solution 1.1}

A faster ending to Solution 1 is as follows.

 
 \begin{align*} f(1)&=(1+p)(1^3+a\cdot1^2+1+10)\\ &=(91)(-77)\\ &= (7)(13)(11)(-7) = (1001)(-7) \\ &=\boxed{\bold{(C)}\, -7007}.\\ \end{align*}


 



\textbf{Solution 1.2}

Also a faster ending to Solution 1 is as follows.

 To find $f(1)$ we just need to find the sum of the coefficients which is $1 + 1 - 8009 + 100 + 900= \boxed{\bold{(C)} \ , -7007}.$

 



\textbf{Solution 2}

We notice that the constant term of $f(x)=c$ and the constant term in $g(x)=10$. Because $f(x)$ can be factored as $g(x) \cdot (x- r)$ (where $r$ is the unshared root of $f(x)$, we see that using the constant term, $-10 \cdot r = c$ and therefore $r = -\frac{c}{10}$.Now we once again write $f(x)$ out in factored form: 

 
 \[f(x) = g(x)\cdot (x-r) = (x^3+ax^2+x+10)(x+\frac{c}{10})\].

 We can expand the expression on the right-hand side to get:

 
 \[f(x) = x^4+(a+\frac{c}{10})x^3+(1+\frac{ac}{10})x^2+(10+\frac{c}{10})x+c\]

 Now we have $f(x) = x^4+(a+\frac{c}{10})x^3+(1+\frac{ac}{10})x^2+(10+\frac{c}{10})x+c=x^4+x^3+bx^2+100x+c$.

 Simply looking at the coefficients for each corresponding term (knowing that they must be equal), we have the equations:
 \[10+\frac{c}{10}=100 \Rightarrow c=900\]
 \[a+\frac{c}{10} = 1, c=900 \Rightarrow a + 90 =1 \Rightarrow a= -89\]

 and finally,

 
 \[1+\frac{ac}{10} = b = 1+\frac{-89 \cdot 900}{10} = b = -8009\].

 We know that $f(1)$ is the sum of its coefficients, hence $1+1+b+100+c$. We substitute the values we obtained for $b$ and $c$ into this expression to get $f(1) = 1 + 1 + (-8009) + 100 + 900 = \boxed{\textbf{(C)}\,-7007}$.

 



\textbf{Solution 3}

Let $r_1,r_2,$ and $r_3$ be the roots of $g(x)$. Let $r_4$ be the additional root of $f(x)$. Then from Vieta's formulas on the quadratic term of $g(x)$ and the cubic term of $f(x)$, we obtain the following:

 
 \begin{align*} r_1+r_2+r_3&=-a \\  r_1+r_2+r_3+r_4&=-1 \end{align*}


 Thus $r_4=a-1$.

 Now applying Vieta's formulas on the constant term of $g(x)$, the linear term of $g(x)$, and the linear term of $f(x)$, we obtain:

 
 \begin{align*} r_1r_2r_3  & = -10\\ r_1r_2+r_2r_3+r_3r_1 &= 1\\  r_1r_2r_3+r_2r_3r_4+r_3r_4r_1+r_4r_1r_2  & = -100\\ \end{align*}


 Substituting for $r_1r_2r_3$ in the bottom equation and factoring the remainder of the expression, we obtain:

 
 \[-10+(r_1r_2+r_2r_3+r_3r_1)r_4=-10+r_4=-100\]

 It follows that $r_4=-90$. But $r_4=a-1$ so $a=-89$

 Now we can factor $f(x)$ in terms of $g(x)$ as

 
 \[f(x)=(x-r_4)g(x)=(x+90)g(x)\]

 Then $f(1)=91g(1)$ and

 
 \[g(1)=1^3-89\cdot 1^2+1+10=-77\]

 Hence $f(1)=91\cdot(-77)=\boxed{\textbf{(C)}\,-7007}$.

 



\textbf{Solution 4 (Risky)}

Let the roots of $g(x)$ be $r_1$, $r_2$, and $r_3$. Let the roots of $f(x)$ be $r_1$, $r_2$, $r_3$, and $r_4$. From Vieta's, we have:
 \begin{align*} r_1+r_2+r_3=-a \\ r_1+r_2+r_3+r_4=-1 \\ r_4=a-1 \end{align*}
The fourth root is $a-1$. Since $r_1$, $r_2$, and $r_3$ are common roots, we have:
 \begin{align*} f(x)=g(x)(x-(a-1)) \\ f(1)=g(1)(1-(a-1)) \\ f(1)=(a+12)(2-a) \\ f(1)=-(a+12)(a-2) \\ \end{align*}
Let $a-2=k$:
 \begin{align*} f(1)=-k(k+14) \end{align*}
Note that $-7007=-1001\cdot(7)=-(7\cdot(11)\cdot(13))\cdot(7)=-91\cdot(77)$This gives us a pretty good guess of $\boxed{\textbf{(C)}\, -7007}$.

 



\textbf{Solution 5}

First off, let's get rid of the $x^4$ term by finding $h(x)=f(x)-xg(x)$. This polynomial consists of the difference of two polynomials with $3$ common factors, so it must also have these factors. The polynomial is $h(x)=(1-a)x^3 + (b-1)x^2 + 90x + c$, and must be equal to $(1-a)g(x)$. Equating the coefficients, we get $3$ equations. We will tackle the situation one equation at a time, starting the $x$ terms. Looking at the coefficients, we get $\dfrac{90}{1-a} = 1$. 
 \[\therefore 90=1-a.\] The solution to the previous is obviously $a=-89$. We can now find $b$ and $c$. $\dfrac{b-1}{1-a} = a$, 
 \[\therefore b-1=a(1-a)=-89*90=-8010\] and $b=-8009$. Finally $\dfrac{c}{1-a} = 10$, 
 \[\therefore c=10(1-a)=10*90=900\] Solving the original problem, $f(1)=1 + 1 + b + 100 + c = 102+b+c=102+900-8009=\boxed{\textbf{(C)}\, -7007}$.

 



\textbf{Solution 6}

Simple polynomial division is a feasible method. Even though we have variables, we can equate terms at the end of the division so that we can cancel terms. Doing the division of $\frac{f(x)}{g(x)}$ eventually brings us the final step $(1-a)x^3 + (b-1)x^2 + 90x + c$ minus $(1-a)x^3 - (a-a^2)x^2 + (1-a)x + 10(1-a)$ after we multiply $f(x)$ by $(1-a)$. Now we equate coefficients of same-degree $x$ terms. This gives us $10(1-a) = c, b-1 = a - a^2, 1-a = 90 \Rightarrow a = -89, c = 900, b = -8009$. We are interested in finding $f(1)$, which equals $1^4 + 1^3 -8009\cdot1^2 + 100\cdot1 + 900 = \boxed{\textbf{(C)}\,-7007}$. ~skyscraper

 



\textbf{Solution 7}

We first note that $f(x) = g(x) \cdot q(x) + r(x)$ where $q$ is the quotient function and $r$ is the remainder function.

 Clearly, $r(x) = 0$ because every single root in $g$ is also in $f$, thus implying $g$ divides $f$.So, we wish to find $f(1) = g(1) \cdot q(1)$.

 Such an expression for $g(1)$ is pretty clean here as we can obtain $g(1) = a + 12$, so we rewrite $f(1) = (a + 12) \cdot q(1)$.Well, now we need to know how $q$ is expressed in order to obtain $q(1)$. This motivates us to long divide to obtain the quotient function. After simple long division $q(x) = x + (1 - a)$. In addition, what is left over, namely $r(x)$, has a constant piece of $a + 89$ (you'll see in a few sentences why we only care about particularly the constant piece).

 Now we can write: $f(1) = (a + 12) \cdot (2 - a)$.

 Now, as we have already established $r(x) = 0$ for ALL $x$ that means $r(0)$ or the constant piece is $0$, so $89 + a = 0$, in which we obtain $a = -89$. We now plug this back into our equation for $f(1)$ to get $(-89 + 12)(2 - (89)) = -77 \cdot 91 = \boxed{\textbf{(C)}\,-7007}$. ~triggod

 



\textbf{Solution 8}

Since $\text{deg(f(x))}=\text{deg(g(x))}+1$, we can write
 \[(f(x)=g(x)\left(x-r \right) \Longrightarrow \left(x^3+ax^2+x+10 \right) \left(x-r \right)=x^4+x^3+bx^2+100x+c.\]Expanding this we get
 \[x^4+\left(a-r \right)x^3+\left(1-ar \right)x^2+\left(10-r \right)x-10r=x^4+x^3+bx^2+100+c.\]From coefficient matching we get the following equations:
 \[a-r=1\]
 \[1-ar=b\]
 \[10-r=100\]
 \[c=-10r.\]Obviously, the third equation results in $r$ being $=90$, and plugging that into the fourth equation we get $c=900$. Now, we can also plug in $r=-90$ to the first equation to get $a=-89$. Finally, plugging $a$ and $r$ into the second equation we find that $b=-8009$. We are asked to find $f(1)$ which is simply $b+c+102 = 7007 \Longrightarrow \boxed{C}$.

 -jb2015007

 



\textbf{Solution 9 (Vieta System)}

Let $p,q,$ and $r$ be the roots of $g(x),$ and $s$ be the additional root of $f(x).$ By Vieta's formulas, it follows that
 \[\begin{cases}p+q+r=-a &(1)\\pq+pr+qr=1 &(2)\\pqr=-10 &(3)\\ p+q+r+s=-1 &(4)\\ pqrs=c &(5)\\pq+pr+ps+qr+qs+rs=b &(6)\\ pqr+pqs+prs+qrs=-100 &(7)\end{cases}\]Looking at $(7),$ we have that
 \[pqr+s(pq+pr+qr)=-10+s=-100\implies s=-90\]then looking at $(5)$ we have that 
 \[c=-10\cdot-90=900\]looking at $(4)$ yields
 \[-1-s=-a\implies a=s+1=-89\]and looking at $(6)$ yields
 \[1+s(p+q+r)=1-as=1-90\cdot89=-8009=b.\]Thus,
 \[f(1)=1+1-8009+100+900=\boxed{\textbf{(C) }-7007}\]~shockfront99

 \textbf{General Notes} $f(1)$ for any polynomial is simply the sum of the coefficients of the polynomial.

 $f(x)/g(x) = x+r$ must have real $r$. Both $f$ and $g$ have all real coefficients, and so odd-degree $g$ must have an odd number of real roots, and even-degree $f$ must have an even number of real roots, so $f$'s single additional root must be real.

 
$77 * 91 = 7 * 11 * 7 * 13$, and $7*11*13=1001$ is a good number sense fact to know. It's interesting because $1001= 10^3+1$ and the 3 nearest primes to $10$ are $7,11,13$.

 


	\item (\textit{AMC 10 10A 2017}) How many integers between $100$ and $999$, inclusive, have the property that some permutation of its digits is a multiple of $11$ between $100$ and $999?$ For example, both $121$ and $211$ have this property.


 $\textbf{(A)}\ 226\qquad\textbf{(B)}\ 243\qquad\textbf{(C)}\ 270\qquad\textbf{(D)}\ 469\qquad\textbf{(E)}\ 486$


 
 \textbf{Answer: }A



\textbf{Solution 1}

There are 81 multiples of 11 between $100$ and $999$ inclusive. Some have digits repeated twice, making 3 permutations.

 Others that have no repeated digits have 6 permutations, but switching the hundreds and units digits also yield a multiple of 11. Switching shows we have overcounted by a factor of 2, so assign  $6 \div 2 = 3$ permutations to each multiple.

 There are now 81*3 = 243 permutations, but we have overcounted*. (At this point, we can already conclude that the solution is 226) Some multiples of 11 have $0$ as a digit. Since $0$ cannot be the digit of the hundreds place, we must subtract a permutation for each. 

 There are 110, 220, 330 ... 990, yielding 9 extra permutations

 Also, there are 209, 308, 407...902, yielding 8 more permutations.

 Now, just subtract these 17 from the total (243) to get 226. $\boxed{\textbf{(A) } 226}$

 
                \begin{itemize}
                    \item If short on time, observe that 226 is the only answer choice less than 243, and therefore is the only feasible answer.
                \end{itemize}


                



\textbf{Solution 2}

We note that we only have to consider multiples of $11$ and see how many valid permutations each has. We can do casework on the number of repeating digits that the multiple of $11$ has:

 $\textbf{Case 1:}$ All three digits are the same. By inspection, we find that there are no multiples of $11$ here.

 $\textbf{Case 2:}$ Two of the digits are the same, and the third is different.

 $\textbf{Case 2a:}$There are $8$ multiples of $11$ without a zero that have this property:$121$, $242$, $363$, $484$, $616$, $737$, $858$, $979$.Each contributes $3$ valid permutations, so there are $8 \cdot 3 = 24$ permutations in this subcase.

 $\textbf{Case 2b:}$There are $9$ multiples of $11$ with a zero that have this property:$110$, $220$, $330$, $440$, $550$, $660$, $770$, $880$, $990$.Each one contributes $2$ valid permutations (the first digit can't be zero), so there are $9 \cdot 2 = 18$ permutations in this subcase.

 $\textbf{Case 3:}$ All the digits are different.Since there are $\frac{990-110}{11}+1 = 81$ multiples of $11$ between $100$ and $999$, there are $81-8-9 = 64$ multiples of $11$ remaining in this case. However, $8$ of them contain a zero, namely $209$, $308$, $407$, $506$, $605$, $704$, $803$, and $902$. Each of those multiples of $11$ contributes $2 \cdot 2=4$ valid permutations, but we overcounted by a factor of $2$; every permutation of $209$, for example, is also a permutation of $902$. Therefore, there are $8 \cdot 4 / 2 = 16$. Therefore, there are $64-8=56$ remaining multiples of $11$ without a $0$ in this case. Each one contributes $3! = 6$ valid permutations, but once again, we overcounted by a factor of $2$ (note that if a number ABC is a multiple of $11$, then so is CBA). Therefore, there are $56 \cdot 6 / 2 = 168$ valid permutations in this subcase.

 Adding up all the permutations from all the cases, we have $24+18+16+168 = \boxed{\textbf{(A) } 226}$.

 



\textbf{Solution 3}

We can first overcount and then subtract.We know that there are $81$ multiples of $11$.

 We can then multiply by $6$ for each permutation of these multiples. (Yet some multiples do not have six distinct permutations.)

 Now divide by $2$, because if a number $abc$ with digits $a$, $b$, and $c$ is a multiple of $11$, then $cba$ is also a multiple of $11$ so we have counted the same permutations twice. 

 Basically, each multiple of $11$ has its own $3$ permutations (say $abc$ has $abc$ $acb$ and $bac$ whereas $cba$ has $cba$ $cab$ and $bca$). We know that each multiple of $11$ has at least $3$ permutations because it cannot have $3$ repeating digits.

 Hence we have $243$ permutations without subtracting for overcounting.Now note that we overcounted cases in which we have $0$'s at the start of each number. So, in theory, we could just answer $A$ and then move on.

 If we want to solve it, then we continue.

 We overcounted cases where the middle digit of the number is $0$ and the last digit is $0$.

 Note that we assigned each multiple of $11$ three permutations.

 The last digit is $0$ gives $9$ possibilities where we overcounted by $1$ permutation for each of $110, 220, ... , 990$.

 The middle digit is $0$ gives $8$ possibilities where we overcount by $1$.$605, 704, 803, 902$ and $506, 407, 308, 209$

 Subtracting $17$ gives $\boxed{\textbf{(A) } 226}$.

 Now, we may ask if there is further overlap (i.e if two of $abc$ and $bac$ and $acb$ were multiples of $11$). Thankfully, using divisibility rules, this can never happen, as taking the divisibility rule mod $11$ and adding, we get that $2a$, $2b$, or $2c$  is congruent to $0\ (mod\ 11)$. Since $a, b, c$ are digits, this can never happen as none of them can equal $11$ and they can't equal $0$ as they are the leading digit of a three-digit number in each of the cases.

 



\textbf{Solution 4 (process of elimination on multiple choices)}

Taken from solution three, we notice that there are a total of $81$ multiples of $11$ between $100$ and $999$, and each of them have at most $6$ permutation (and thus is a permutations of $6$ numbers), giving us a maximum of $486$ valid numbers. 

 However, if $abc$ can be divided by $11$, so can $cba$, which is distinct if $c \neq a$. And if $c = a$ then $abc$ and $cba$ have the same permutations. Either way, we have doubled counted. This reduces the number of permutations to $486/2 =  243$. 

 Furthermore, if $a=b$ or $c=0$ (which turn out to be equivalent conditions! for example $220$), not all (inverse) permutations are distinct ($\mathbf{2}20 = 2\mathbf{2}0$) or valid ($022$). (There are 9 of these.) 

 Similarly, for $a0b$, not all (inverse) permutations are valid. (There are 8 of these.)

 As long as you notice at least one example of one of these 3 cases, you may infer that the answer must be smaller than $243$. This leaves us with only one possible answer: $\boxed{\textbf{(A) } 226}$.

 


\end{enumerate}